\section{Introduction}\label{sec:intro}

Reinforcement learning in adversarial settings poses unique challenges where an agent 
must not only learn to act effectively, but to do so against an opponent that actively 
counters its behavior. In this work, we train agents for a 1v1 hockey environment 
inspired by air hockey, where two players compete on opposite sides of a rink to 
score goals against each other. The game requires a combination of offensive and 
defensive skills across a fully continuous action space, with scoring goals as 
the ultimate objective.

\medskip
Generalizing across different opponent play styles proves to be the central difficulty. 
We address this through reward shaping to accelerate early learning, exploration 
design, and most importantly a self-play curriculum that progressively challenges the 
agent with a diverse and evolving pool of opponents. We find that the diversity of 
this pool, rather than any single algorithmic choice, is the primary driver of 
agent robustness.

\medskip
We implement and compare two established off-policy algorithms, \ac{TD3} and \ac{SAC}, 
and analyze their behavior under different training conditions. The algorithms are first 
validated on standard continuous control benchmarks before being applied to the hockey 
environment. Methodology is covered in Section~\ref{sec:method}, experimental results 
in Section~\ref{sec:eval}, and conclusions in Section~\ref{sec:discussion}. 
The full code and configuration files are available at 
\url{https://github.com/JannikRom/drai-rl}.